\centerline{\bf \Huge Геометрия 3}
\centerline{\bf \Huge Касания и углы}

\begin{flushright}
        \scriptsize{Авраам Линкольн дал людям свободу,\\ а полковник Кольт уравнял их шансы}
\end{flushright}

\problem{1} В остроугольном треугольнике $A B C$ проведены высоты $B B_1$ и $C C_1$. Через точку $A$ проводится прямая $\ell$, параллельная прямой $B_1 C_1$. Докажите, что эта прямая касается окружности ( $A B C$ ).

\problem{2} Прямая $P A$ касается описанной окружности треугольника $A B C$. Точки $C_1$ и $B_1$ - основания перпендикуляров, опущенных из $P$ на прямые $A B, A C$. Докажите, что $B C \perp B_1 C_1$.

\problem{3} В остроугольном треугольнике $A B C$ проведены высоты $B B_1$ и $C C_1$. Точка $M$ - середина отрезка $B C$. Докажите, что прямые $M B_1$ и $M C_1$ касаются окружности $\left(B_1 A C_1\right)$.

\problem{4} (a) Касательная к описанной окружности треугольника $A B C$, проведенная через точку $A$, пересекает прямую $B C$ в точке $X$. Биссектриса угла $A$ пересекает сторону $B C$ в точке $L$. Докажите, что треугольник $A L X$ равнобедренный. (b) Сформулируйте и докажите аналогичное утверждение для биссектрисы внешнего угла $A$.

\problem{5} Треугольник $A B C$ вписан окружность $\omega$. Прямая $\ell$ проходит через точку $A$, касается $\omega$ и пересекает луч $B C$ в точке $S$. Обозначим за $X$ проекцию точки $A$ на прямую $B C$, а за $Y$ проекцию точки $S$ на прямую $A B$. Докажите, что $X Y$ делит отрезок $A C$ пополам.

\problem{6} Прямая $\ell$ касается окружности, описанной около треугольника $A B C$ в точке $A$. На прямой $\ell$ выбирается точка $X$, а на отрезке $B C$ точка $Y$ так, что $C X \perp B C$ и $X Y \| A B$. Докажите, что прямая $A Y$ проходит через центр окружности, описанной около треугольника $A B C$.

\problem{7} В неравнобедренном треугольнике $A B C$ проведена высота из вершины $A$ и биссектрисы из двух других вершин. Докажите, что описанная окружность треугольника, образованного этими тремя прямыми, касается биссектрисы, проведенной из вершины $A$.

\newpage

\hspace{3mm}

\problem{8} Остроугольный равнобедренный треугольник $A B C(A B=A C)$ вписан в окружность с центром в точке $O$. Лучи $B O$ и $C O$ пересекают стороны $A C$ и $A B$ в точках $B^{\prime}$ и $C^{\prime}$ соответственно. Через точку $C^{\prime}$ проведена прямая $\ell$, параллельная прямой $A C$. Докажите, что прямая $\ell$ касается окружности, описанной около треугольника $B^{\prime} O C$.

\problem{9$^{*}$}
Точка $E$ выбирается на продолжении стороны $A D$ прямоугольника $A B C D$ за точку $D$. Луч $E C$ пересекает окружность $\omega$, описанную около треугольника $A B E$ в точке $F \neq E$. Лучи $D C$ и $A F$ пересекаются в $P$. Точка $H$ основание перпендикуляра, опущенного из точки $C$ на прямую $\ell$, проходящую через $E$ параллельно $A F$. Докажите, что прямая $P H$ касается $\omega$.

\problem{10$^{*}$}
Точка $P$ выбрана на меньшей дуге $B C$ окружности $\Omega$, описанной около остроугольного треугольника $A B C$. Точка $Q$ выбрана внутри треугольника $A B C$ так, что $\angle P B Q=\angle A C B$ и $\angle P C Q=\angle A B C$. Точка $R \neq Q$ выбрана на отрезке $A P$ так, что $Q R \| B C$. Наконец, точка $S \neq A$ выбрана на дуге $B P$ окружности $\Omega$ так, что $R A=R S$. Докажите, что точки $P, Q, R$ и $S$ лежат на одной окружности.